\documentclass[11pt]{article}
\usepackage[margin=1in]{geometry}
\usepackage{amsmath,amssymb}
\usepackage{hyperref}

\title{G8 Lane A A1/A2 Operator Self-Adjoint Proof Map}
\author{Private red-team proof-map memo}
\date{May 21, 2026}

\newcommand{\Spec}{\operatorname{Spec}}
\newcommand{\Dom}{\operatorname{Dom}}
\newcommand{\im}{\operatorname{Im}}
\newcommand{\Mem}{\operatorname{Mem}}
\newcommand{\ptspec}{\operatorname{Spec}_{p}}

\begin{document}
\maketitle

\section*{Executive Verdict}

Lane A now has three explicit axiom gates.  This memo concerns only the first
two:
\[
\begin{aligned}
&\text{A1: construct a ready Book III lemniscate operator package},\\
&\text{A2: construct an operator-native self-adjoint spectral predicate}.
\end{aligned}
\]
The current Lean spine correctly treats A1 as a genuine operator-readiness
theorem.  It is not finite-check plumbing.  The existing finite
\(n^2\)-readout and finite self-adjointness checks are useful diagnostics, but
they do not by themselves construct the analytic graph operator or its ready
self-adjoint spectral package.

The A2 route should be the analytic point spectrum of the actual Chapter~23
lemniscate graph operator.  The standard finite/typed \(n^2\) spectrum remains
compatibility evidence unless it is separately proved to coincide with that
operator point spectrum.

\section*{Lean Targets}

The Lane-A ledger names the two gates as:
\begin{itemize}
\item \texttt{G8LaneAOperatorReadyGate}.
\item \texttt{G8LaneAOperatorNativeSelfAdjointGate}.
\end{itemize}
Expanded, A1 is:
\begin{itemize}
\item \texttt{operatorCtx : LemniscateOperatorContext}.
\item \texttt{operatorReady : LemniscateOperatorReady operatorCtx}.
\end{itemize}
The readiness predicate requires:
\begin{itemize}
\item \texttt{LemniscateDomainReady operatorCtx.domain}.
\item \texttt{LemniscateSelfAdjointObligation operatorCtx}.
\item \texttt{LemniscateCompactResolventObligation operatorCtx}.
\item \texttt{LemniscateDiscreteSpectrumObligation operatorCtx}.
\item \texttt{operatorCtx.obligations.obligationsRespectDomain}.
\item \texttt{operatorCtx.finiteApproximationCompatible}.
\end{itemize}
A2 should be discharged by a strict provenance package:
\[
  \texttt{G8OperatorNativeSpectralProvenance operatorCtx operatorReady}.
\]
Its load-bearing fields are:
\begin{itemize}
\item \texttt{spectralMembership}: the exported complex spectral predicate for
  the ready operator package.
\item \texttt{nativeMembership : C -> Prop}.
\item \texttt{kind : G8OperatorNativeSpectralProvenanceKind}.
\item \texttt{spectralMembership\_iff\_native}.
\item \texttt{selfAdjointReality\_native}: every native member has
  \texttt{lambda.im = 0}.
\end{itemize}
The final A1/A2 package is:
\[
  \texttt{G8LaneAOperatorNativeSelfAdjointSource}.
\]
It must not mention actual \(\xi\)-carriers.  Canonical actual-\(\xi\)
membership is A3.

\section*{Source Anchors}

The source truth for this proof map is Book III, Chapter~23:
\begin{itemize}
\item Ch.~23, Definition III.D28: \(H_L=-d^2/dx^2\) on the two-lobe
  lemniscate metric graph with continuity and Kirchhoff boundary conditions.
\item Ch.~23, Theorem III.T17: \(H_L\) is self-adjoint by the boundary-form
  computation and maximality of Kirchhoff boundary conditions.
\item Ch.~23, Proposition III.P09: compactness of the metric graph gives
  compact resolvent and discrete real spectrum.
\item Ch.~23, spectral theorem discussion: the self-adjoint operator with
  compact resolvent has point spectrum and spectral resolution.
\end{itemize}
Chapters~24--25 explain why this operator is useful downstream, but they are
not A1/A2 proof sources.  In particular, determinant correspondence,
actual-\(\xi\) membership, critical-line forcing, and accepted tower coverage
belong to later gates.

\section*{Mathematical Proof Map}

\paragraph{Step 1: construct the metric graph.}
Let
\[
  L=S^1_B\vee S^1_C
\]
be the lemniscate graph with two circle lobes meeting at the crossing point
\(\omega\).  Each lobe is parametrized by an interval with endpoints
identified at \(\omega\).  This supplies the geometric carrier for the
operator context.

\paragraph{Step 2: construct the Hilbert space and domain.}
Define
\[
  L^2(L)=L^2(S^1_B)\oplus L^2(S^1_C).
\]
The operator domain is the Sobolev/Kirchhoff domain
\[
  \Dom(H_L)=
  \left\{
    f=(f_B,f_C):
    f_B,f_C\in H^2,\
    f_B(\omega)=f_C(\omega),\
    f'_B(\omega)+f'_C(\omega)=0
  \right\}.
\]
In Lean-facing terms, this step must prove the fields of
\texttt{LemniscateDomainReady}: Hilbert readiness, value-trace definition and
continuity, derivative-trace definition and continuity, crossing closure, and
Kirchhoff closure.

\paragraph{Step 3: define the operator.}
Define
\[
  H_L f = -\frac{d^2 f}{dx^2}
\]
edgewise on \(S^1_B\) and \(S^1_C\), with domain \(\Dom(H_L)\).  This
replaces the current scaffold placeholder by the analytic graph Laplacian.

\paragraph{Step 4: prove symmetry by boundary-form cancellation.}
For \(f,g\in\Dom(H_L)\), integration by parts on both lobes gives
\[
  \langle H_Lf,g\rangle-\langle f,H_Lg\rangle
  =
  -\overline{g(\omega)}(f'_B(\omega)+f'_C(\omega))
  +f(\omega)\overline{(g'_B(\omega)+g'_C(\omega))}.
\]
The Kirchhoff conditions for \(f\) and \(g\) make both sums vanish.  Hence the
boundary form is zero and \(H_L\) is symmetric.

\paragraph{Step 5: prove self-adjointness.}
The Kirchhoff boundary condition is maximal among boundary conditions yielding
a symmetric extension for the one-dimensional graph Laplacian.  Equivalently,
standard quantum-graph theory gives that the Laplacian on a compact metric
graph with Kirchhoff vertex conditions is self-adjoint.  This proves:
\[
  H_L^*=H_L.
\]
This is the load-bearing A1 self-adjointness field.  It cannot be replaced by
finite checks.

\paragraph{Step 6: prove compact resolvent and discrete spectrum.}
The graph \(L\) is compact.  The Sobolev embedding from the operator domain
into \(L^2(L)\) is compact, so the resolvent \((H_L-\lambda)^{-1}\) is compact
for \(\lambda\) off the spectrum.  Therefore \(H_L\) has discrete point
spectrum with finite-multiplicity eigenvalues and no continuous-spectrum
residue relevant to this Lane-A predicate.

\paragraph{Step 7: package A1.}
Build a `LemniscateOperatorContext` whose obligations are the theorems proved
above, and prove:
\[
  \texttt{LemniscateOperatorReady operatorCtx}.
\]
Then construct:
\[
  \texttt{G8LaneAOperatorReadyGate}.
\]

\paragraph{Step 8: define the analytic point-spectrum predicate.}
For a ready operator context, define the native spectral predicate
\[
  \Mem_{H_L}(\lambda)
  \quad:\Longleftrightarrow\quad
  \lambda\in\ptspec(H_L)
\]
where the real eigenvalue is read as a complex value \(r+0i\).  Lean should use
this predicate as both the exported `spectralMembership` and the
`nativeMembership`, or prove exact equivalence if the exported wrapper differs.
Set:
\[
  \texttt{kind := .pointSpectrum}.
\]

\paragraph{Step 9: prove self-adjoint spectral reality.}
If \(\lambda\in\ptspec(H_L)\), choose a nonzero eigenvector \(\psi\) with
\[
  H_L\psi=\lambda\psi.
\]
By self-adjointness,
\[
  \lambda\|\psi\|^2
  =
  \langle H_L\psi,\psi\rangle
  =
  \langle \psi,H_L\psi\rangle
  =
  \overline{\lambda}\|\psi\|^2.
\]
Since \(\psi\ne0\), \(\|\psi\|^2\ne0\), so \(\lambda=\overline{\lambda}\), and
therefore
\[
  \im(\lambda)=0.
\]
This is the load-bearing A2 theorem:
\[
  \forall \lambda,\ \Mem_{H_L}(\lambda)\to\lambda.\operatorname{im}=0.
\]

\paragraph{Step 10: package A2 and the combined source.}
Construct:
\[
  \texttt{G8OperatorNativeSpectralProvenance operatorCtx operatorReady}
\]
with:
\begin{itemize}
\item \texttt{spectralMembership := analyticPointSpectrum}.
\item \texttt{nativeMembership := analyticPointSpectrum}.
\item \texttt{kind := .pointSpectrum}.
\item \texttt{spectralMembership\_iff\_native := fun \_ => Iff.rfl}.
\item \texttt{selfAdjointReality\_native := pointSpectrum\_real}.
\end{itemize}
Then construct:
\[
  \texttt{G8LaneAOperatorNativeSelfAdjointSource}.
\]
Existing adapters then produce the Lane-A ready gate and operator-native
self-adjoint gate.  They do not reduce the A3 membership gate.

\section*{Lean Follow-Up Targets}

The next Lean wave should first split a pure A1/A2 core with no actual-\(\xi\)
imports.  Recommended theorem-facing names:
\[
\begin{gathered}
  \texttt{G8BookIIILemniscateOperatorReadySource},\\
  \texttt{G8BookIIILemniscateAnalyticPointSpectrumProvenance},\\
  \texttt{G8BookIIILaneAOperatorNativeSelfAdjointSource}.
\end{gathered}
\]
The clean A1/A2 module should import only operator/core spectral definitions
and must avoid final-spine or actual-\(\xi\) membership modules.

\section*{Guardrails}

\begin{itemize}
\item Do not use A3 canonical actual-\(\xi\) membership to construct A1 or A2.
\item Do not use determinant transfer, O3, completion uniqueness, divisor
  transfer, accepted coverage, final live hinge, or Mathlib discharge modules.
\item Do not promote finite checks such as \texttt{self\_adjoint\_check 20},
  \texttt{k5\_diagonal\_check 20}, or finite spectral-resolution checks into global
  operator readiness.
\item Do not treat the standard \(n^2\) readout as the analytic operator point
  spectrum unless equality with the actual Chapter~23 operator spectrum is
  separately proved.
\item Keep ledger axiom counts unchanged until a theorem-backed constructor is
  actually available.
\end{itemize}

\section*{Decision Point}

The mathematical route for A1/A2 is clear and non-circular:
\[
\begin{aligned}
&\text{compact Kirchhoff metric graph operator}\\
&\quad\Rightarrow\text{ self-adjoint ready operator}\\
&\quad\Rightarrow\text{ real analytic point spectrum}\\
&\quad\Rightarrow\text{ strict A1/A2 source}.
\end{aligned}
\]
The remaining work is not bridge plumbing.  It is the Lean formalization of the
Chapter~23 compact-metric-graph operator theorem in exactly the form consumed
by \texttt{LemniscateOperatorReady} and
\texttt{G8OperatorNativeSpectralProvenance}.

\end{document}
